\chapter{单片机脚本运行功能} \label{ch:script-engine}

“脚本”通常指不需要重新编译、只修改文本就能改变程序行为的机制。单片机上的脚本不要求像 Python 那样完整，最常见的需求是：用户通过串口写入一段表达式，设备能解析并计算出结果。本章从一个计算表达式的小解释器入手，介绍词法、语法、递归下降解析、错误处理与堆栈保护，再简单讨论 $LL(1)$ 文法与 Lua 框架。完整示例代码位于 \texttt{Chapter-script-engine/表达式脚本引擎/} 目录。

\medskip 本章涉及到的实验内容：
\begin{enumerate}
\item 表达式词法分析（数字、变量和运算符）；
\item 递归下降语法分析；
\item 语法错误、除零、未知变量处理；
\item 递归深度限制与堆栈溢出防护。
\end{enumerate}

% ---------------------------------------------------------------------
\section{脚本功能}

在资源有限的单片机上，脚本引擎并不需要“万能”。本章的目标是让设备支持类似下面的输入：

\begin{lstlisting}
10 + 2 * (3 + x)
(a + b) * 2
100 / 4 - 3 * 2
\end{lstlisting}

它的价值在于：

\begin{itemize}
\item 参数可在线修改：把“阈值、系数、限位”交给用户配置，不用为每个参数写一个特殊协议；
\item 逻辑与代码分离：公式可以放在配置区、Flash 或上位机下发，固件只维护一个解释器；
\item 表达式比固定格式更灵活：固定协议需要设计很多命令，表达式只需一条 \texttt{calc} 命令。
\end{itemize}

对应的运行时上下文只需要一个变量表：

\begin{lstlisting}[style=c]
typedef struct __ExprVar_T
{
    char name[EXPR_NAME_MAX]; /* 变量名 */
    int32_t value;            /* 变量值 */
    uint8_t used;             /* 是否已使用 */
} ExprVar_T;

typedef struct __ExprContext_T
{
    ExprVar_T vars[EXPR_VAR_MAX];
    uint16_t varCount;
} ExprContext_T;
\end{lstlisting}

% ---------------------------------------------------------------------
\section{语法分析器}

\subsection{什么是文法}

“文法”（Grammar）是用一组规则描述“什么样的符号序列才是合法程序或表达式”的形式化定义。自然语言有语法，编程语言同样有语法。例如表达式 \texttt{10 + 2 * (3 + x)} 合法，而 \texttt{10 + * 2} 不合法——判断依据不是“能不能算出数”，而是它是否满足表达式的文法。

文法通常由四部分组成：

\begin{itemize}
\item 终结符（Terminal）：不能再展开的最小单元，例如数字、变量名、 \texttt{+ - * / ( )}；
\item 非终结符（Nonterminal）：可以继续展开的语法成分，例如 \texttt{expr}、\texttt{term}、\texttt{factor}；
\item 产生式（Production）：描述“非终结符可以展开成什么”，例如 \texttt{expr ::= term ('+' term)*}；
\item 开始符号（Start Symbol）：整个语言从哪个非终结符开始推导。
\end{itemize}

如果把文法比作“语法地图”，词法分析负责把字符拆成单词，语法分析则负责按这张地图判断单词的排列是否合法、结构是什么。

\subsection{为什么要进行语法分析}

语法分析要回答两个问题：

\begin{enumerate}
\item 输入是否符合文法：表达式是否完整、括号是否匹配、运算符位置是否正确；
\item 如果符合，结构是什么：哪部分是加法、哪部分是乘法、哪部分要先算。
\end{enumerate}

没有语法分析，程序只能收到“一串字符”，无法知道 \texttt{2 * (3 + x)} 是一个整体；有了语法分析，解析器才能按照优先级与结合性生成结构，再交给求值器、AST 和虚拟机执行。它同时是错误报告的第一道关口：语法错误应该在“结构层”被拦截，而不是等到运行时算出一个奇怪的结果。

语法分析前必须先做词法分析（Lexer，也称 Scanner）：把字符串转成有意义的最小单元。本章示例支持数字、变量名、四则运算符和括号，Token 定义如下：

\begin{lstlisting}[style=c]
typedef enum __ExprTokenKind_T
{
    EXPR_TOKEN_NONE = 0,
    EXPR_TOKEN_NUMBER,
    EXPR_TOKEN_IDENT,
    EXPR_TOKEN_PLUS,
    EXPR_TOKEN_MINUS,
    EXPR_TOKEN_STAR,
    EXPR_TOKEN_SLASH,
    EXPR_TOKEN_LPAREN,
    EXPR_TOKEN_RPAREN,
    EXPR_TOKEN_END
} ExprTokenKind_T;
\end{lstlisting}

扫描器一次只读取一个“向前看”的 Token（Lookahead），语法层通过这个 Token 决定下一步动作。这种“扫描器 + 当前 Token”的结构是所有解析器的通用骨架。

语法分析器依据文法判断 Token 序列是否合法。本章表达式的文法如下：

\begin{lstlisting}
expr   := term (('+' | '-') term)*
term   := factor (('*' | '/') factor)*
factor := number | ident | '(' expr ')' | ('-' | '+') factor
\end{lstlisting}

文法分层的原因是运算符优先级：\texttt{*} 与 \texttt{/} 的优先级高于 \texttt{+} 与 \texttt{-}，因此乘法所在的 \texttt{term} 必须嵌套在加法所在的 \texttt{expr} 内部。

% ---------------------------------------------------------------------
\section{递归下降器的实现}

递归下降（Recursive Descent）让每一个语法成分对应一个函数：

\subsection{实现原理}

递归下降解析器的核心思想只有三条：

\begin{enumerate}
\item 一个非终结符对应一个函数：文法的 \texttt{expr}、\texttt{term} 与 \texttt{factor} 分别对应 \texttt{Expr\_ParseExpr}、\texttt{Expr\_ParseTerm} 和 \texttt{Expr\_ParseFactor}；
\item 用“当前 Lookahead Token”决定走哪个分支：函数先看扫描器里的当前 Token，判断接下来该匹配数字、变量、括号还是运算符；
\item 递归调用代替“压栈和弹栈”：遇到括号里的表达式时， \texttt{Expr\_ParseFactor} 调用 \texttt{Expr\_ParseExpr}，新表达式的解析状态由新的函数调用栈保存，解析完括号再返回上一层继续。
\end{enumerate}

优先级并不是由运算符优先级表实现的，而是靠文法分层实现的：

\begin{itemize}
\item \texttt{expr} 先调用 \texttt{term}，因此加减法比乘除法“外层”；
\item \texttt{term} 再调用 \texttt{factor}，因此乘除法比括号和变量更“外层”；
\item 只有解析完一个完整的 \texttt{term}、\texttt{factor}，才会回到 \texttt{expr} 判断是否有下一个 \texttt{+} 或 \texttt{-}。
\end{itemize}

递归下降不需要显式构造语法树：本章示例采用“边解析边求值”，每个函数解析完子表达式后，立刻把结果写入 \texttt{pValue} 返回给上一层。若脚本语言需要先编译后执行，则只需把“写 \texttt{pValue}”改成“创建 AST 节点并返回”。

\begin{lstlisting}[style=c]
/* expr := term (('+'|'-') term)* */
static int8_t Expr_ParseExpr(ExprContext_T *pCtx,
                             ExprScanner_T *pScanner,
                             int32_t *pValue)
{
    int8_t ret = Expr_ParseTerm(pCtx, pScanner, pValue);

    while ((ret == EXPR_ERR_OK) &&
           ((pScanner->token == EXPR_TOKEN_PLUS) ||
            (pScanner->token == EXPR_TOKEN_MINUS))) {
        ExprTokenKind_T op = pScanner->token;

        ret = Expr_ScanNext(pScanner);
        if (ret != EXPR_ERR_OK) {
            break;
        }

        int32_t right = 0;
        ret = Expr_ParseTerm(pCtx, pScanner, &right);
        if (op == EXPR_TOKEN_PLUS) {
            *pValue = *pValue + right;
        } else {
            *pValue = *pValue - right;
        }
    }
    return ret;
}
\end{lstlisting}

\texttt{Expr\_ParseFactor} 处理数字、变量、一元正负号和括号：

\begin{lstlisting}[style=c]
static int8_t Expr_ParseFactor(ExprContext_T *pCtx,
                               ExprScanner_T *pScanner,
                               int32_t *pValue)
{
    switch (pScanner->token) {
        case EXPR_TOKEN_NUMBER:
            *pValue = (int32_t)strtol(pScanner->tokenText, NULL, 10);
            return Expr_ScanNext(pScanner);

        case EXPR_TOKEN_IDENT:
            /* 查变量表 */
            ...

        case EXPR_TOKEN_LPAREN:
            ret = Expr_ScanNext(pScanner);
            ret = Expr_ParseExpr(pCtx, pScanner, pValue);
            if (pScanner->token != EXPR_TOKEN_RPAREN) {
                return EXPR_ERR_SYNTAX;   /* 括号不匹配 */
            }
            return Expr_ScanNext(pScanner);

        default:
            return EXPR_ERR_SYNTAX;
    }
}
\end{lstlisting}

\subsection{运行流程}

以表达式 \texttt{10 + 2 * (3 + x)} 为例，假设 \texttt{x = 10}，递归下降的执行过程如下：

\begin{figure}[htbp]
\centering
	\begin{tikzpicture}[
		scale=0.7,            % 新增，整体缩放
		nt/.style={draw=structurecolor, thick, rounded corners=2pt,
			fill=structurecolor!8, align=center, font=\small,
			minimum width=2.0cm, minimum height=0.8cm},
		tm/.style={draw=structurecolor!60, thick, rounded corners=2pt,
			fill=white, align=center, font=\small,
			minimum width=1.5cm, minimum height=0.7cm},
		edge/.style={-stealth, thick, structurecolor}
		]
		\node[nt] (expr)  at (0,0)    {expr};
		\node[nt] (term1) at (-3,-2.2) {term};
		\node[tm] (plus)  at (0,-2.2) {$+$};
		\node[nt] (term2) at (3.2,-2.2) {term};
		
		\node[nt] (f1) at (-3,-4.4) {factor};
		\node[tm] (v10) at (-3,-6.2) {10};
		
		\node[nt] (f2) at (0.3,-4.4) {factor};
		\node[tm] (v2) at (0.3,-6.2) {2};
		\node[tm] (mul) at (3.2,-4.4) {$*$};
		\node[nt] (fp) at (6.1,-4.4) {factor};
		\node[nt] (expr2) at (6.1,-6.6) {expr};
		
		\node[nt] (term3) at (3.4,-8.6) {term};
		\node[tm] (plus2) at (6.1,-8.6) {$+$};
		\node[nt] (term4) at (8.8,-8.6) {term};
		\node[nt] (f3) at (3.4,-10.6) {factor};
		\node[tm] (v3) at (3.4,-12.3) {3};
		\node[nt] (f4) at (8.8,-10.6) {factor};
		\node[tm] (vx) at (8.8,-12.3) {x = 10};
		
		\draw[edge] (expr.south) -- (term1.north);
		\draw[edge] (expr.south) -- (plus.north);
		\draw[edge] (expr.south) -- (term2.north);
		\draw[edge] (term1.south) -- (f1.north);
		\draw[edge] (f1.south) -- (v10.north);
		\draw[edge] (term2.south) -- (f2.north);
		\draw[edge] (term2.south) -- (mul.north);
		\draw[edge] (term2.south) -- (fp.north);
		\draw[edge] (f2.south) -- (v2.north);
		\draw[edge] (fp.south) -- (expr2.north);
		\draw[edge] (expr2.south) -- (term3.north);
		\draw[edge] (expr2.south) -- (plus2.north);
		\draw[edge] (expr2.south) -- (term4.north);
		\draw[edge] (term3.south) -- (f3.north);
		\draw[edge] (f3.south) -- (v3.north);
		\draw[edge] (term4.south) -- (f4.north);
		\draw[edge] (f4.south) -- (vx.north);
		
		\node[font=\small, align=left] at (7.8,-2.2) {运算符为终结符，\\不需要继续展开};
		\node[font=\small, align=left] at (8.8,-6.6) {括号触发\\递归 expr};
	\end{tikzpicture}
\caption{\texttt{10 + 2 * (3 + x)} 的递归下降解析树} \label{fig:recursive-descent-tree}
\end{figure}

\begin{enumerate}
\item \texttt{Expr\_ParseExpr} 调用 \texttt{Expr\_ParseTerm}；
\item \texttt{Expr\_ParseTerm} 调用 \texttt{Expr\_ParseFactor}，读数字 10，继续读到 \texttt{'+'}，因此 term 结束，返回 10；
\item \texttt{Expr\_ParseExpr} 看到 \texttt{'+'}，消费加号，再调用一次 \texttt{Expr\_ParseTerm} 解析右操作数；
\item 新的 term 先解析出数字 2，看到 \texttt{'*'} 后继续调用 factor 解析 \texttt{(3 + x)}；
\item 括号触发 \texttt{Expr\_ParseFactor} 递归进入新的 \texttt{Expr\_ParseExpr}，内部算出 \texttt{3 + 10 = 13}；
\item 回到乘法层：\texttt{2 * 13 = 26}；
\item 回到加法层：\texttt{10 + 26 = 36}，并确认输入已结束。
\end{enumerate}

整个过程的调用顺序如表~\ref{tab:recursive-flow}~所示。

\begin{table}[htbp]
\centering \caption{\texttt{10 + 2 * (3 + x)} 的递归下降运行流程} \label{tab:recursive-flow}
    \begin{tabular}{lll}
        \toprule
        步骤 & 动作 & 当前结果 \\
        \midrule
        1 & expr -> term -> factor 读取 10 & 10 \\
        2 & 遇到 \texttt{'+'}，准备解析右项 & 10 \\
        3 & term -> factor 读取 2 & 2 \\
        4 & 遇到 \texttt{'*'}，factor 读取 \texttt{'('} & 2 \\
        5 & 括号内 expr：3 + x & 13 \\
        6 & 括号结束，乘法：2 * 13 & 26 \\
        7 & 加法：10 + 26 & 36 \\
        \bottomrule
    \end{tabular}
\end{table}

从运行流程可以看出递归下降的两个特点：每层函数“只负责自己那一条语法规则”；括号越深，函数调用栈越深，这也为后续“递归深度限制”提供了理论依据。

递归下降的优点是代码结构与文法一一对应、容易阅读；缺点是函数之间互相递归，编译器必须允许一定深度的调用栈，这也是本章把“防止堆栈溢出”单独拿出来讨论的原因。

% ---------------------------------------------------------------------
\subsection{错误处理的实现}

解释器不能像编译器那样把错误全部报完再退出，而应采用“快速失败”（Fail Fast）策略：遇到第一个错误立即返回错误码，停止后续解析。本章定义的错误码包括：

\begin{lstlisting}
EXPR_ERR_SYNTAX       语法错误（非法字符、括号不匹配、表达式不完整）
EXPR_ERR_UNKNOWN_VAR  引用了未设置的变量
EXPR_ERR_DIV_ZERO     除数为零
EXPR_ERR_DEPTH_LIMIT  括号或递归嵌套超过上限
EXPR_ERR_VAR_FULL     变量表已满
\end{lstlisting}

公共入口 \texttt{Expr\_Evaluate} 在返回前必须保证：

\begin{enumerate}
\item 参数非空；
\item 词法扫描没有非法字符；
\item 语法函数返回成功；
\item 整条表达式都被消费完，最后一个 Token 必须是 END。
\end{enumerate}

最后一条很关键。例如 \texttt{10 + (2)} 合法，但 \texttt{10 + (2))} 会留下多余的右括号，在末尾检查 \texttt{scanner.token != EXPR\_TOKEN\_END} 时被拒绝。

% ---------------------------------------------------------------------
\subsection{如何防止堆栈溢出}

递归下降解析器的栈用量与“表达式嵌套深度”成正比：

\begin{lstlisting}
1 + (2 + (3 + (4 + ... )))
\end{lstlisting}

括号嵌套越深，\texttt{ParseFactor} 与 \texttt{ParseExpr} 的递归层数越多；如果脚本由上位机任意下发，攻击者或误操作可能让设备栈溢出。本章采用两层保护：

\begin{enumerate}
\item 递归深度限制：扫描器中记录 \texttt{depth}，每进入一层 factor 就加 1，超过 \texttt{EXPR\_MAX\_PARSE\_DEPTH} 立即返回错误；
\item MCU 栈规划：解析器所在的调用栈必须预留足够空间，不能放在只有几百字节的中断上下文里。
\end{enumerate}

\begin{lstlisting}[style=c]
if (pScanner->depth >= EXPR_MAX_PARSE_DEPTH) {
    return EXPR_ERR_DEPTH_LIMIT;
}
pScanner->depth++;

/* ... 解析当前 factor ... */

pScanner->depth--;
return ret;
\end{lstlisting}

深度限制必须与 MCU 栈大小一起评估：限制设得太大，函数本身就会把栈占满；限制设得太小，合法表达式反而无法解析。

% ---------------------------------------------------------------------
\section{LL(1) 文法}

$LL(1)$ 是“自顶向下、从左到右扫描、最左推导、只需向前看 1 个 Token”的语法类别。本章文法满足 $LL(1)$ 的基本要求：

\begin{enumerate}
\item 每个产生式右侧的第一个符号不同，或能通过一个 lookahead 区分；
\item 不包含左递归。例如不能写成 \texttt{expr ::= expr '+' term}，因为解析器会无限调用 \texttt{ParseExpr}；
\item 不产生二义性。
\end{enumerate}

\begin{anybox}{左递归与生成器}
这里“不允许左递归”针对的是手工递归下降解析器；本章后面的 .gram 生成器会把
\texttt{expr ::= expr PLUS term} 自动改写成“基产生式 + 循环续产生式”，
因此可以直接使用左递归写法。
\end{anybox}

\subsection{First 与 Follow 集}

判断“当前 Lookahead Token 应该走哪个产生式”，需要两个形式化工具：

\begin{itemize}
\item $\mathrm{First}(\alpha)$：从符号串 $\alpha$ 出发，所有可能“第一个读到的终结符”集合；若 $\alpha$ 能推导出空串(ε)，则 $\varepsilon \in \mathrm{First}(\alpha)$；
\item $\mathrm{Follow}(A)$：在某个句型中，紧跟非终结符 $A$ 之后可能出现的终结符集合；若 $A$ 可能出现在句型末尾，则 \texttt{\$}（输入结束符）也属于 $\mathrm{Follow}(A)$。
\end{itemize}

以无左递归的表达式文法为例：

\begin{lstlisting}
expr  -> term expr'
expr' -> '+' term expr' | ε
term  -> factor term'
term' -> '*' factor term' | ε
factor-> '(' expr ')' | number
\end{lstlisting}

各非终结符的 First 与 Follow 集合如表~\ref{tab:first-follow}~所示。

\begin{table}[htbp]
\centering \caption{表达式文法的 First 与 Follow 集合} \label{tab:first-follow}
    \begin{tabular}{llll}
        \toprule
        非终结符 & First & Follow & 含义 \\
        \midrule
        \texttt{expr}  & $\{number, (\}$ & $\{\$, )\}$ & 表达式可以从数字或括号开始 \\
        \texttt{expr'} & $\{+, \varepsilon\}$ & $\{\$, )\}$ & 加号后缀可空 \\
        \texttt{term}  & $\{number, (\}$ & $\{+, \$, )\}$ & term 结束后可能遇到加号 \\
        \texttt{term'} & $\{*, \varepsilon\}$ & $\{+, \$, )\}$ & 乘号后缀可空 \\
        \texttt{factor}& $\{number, (\}$ & $\{*, +, \$, )\}$ & factor 后可继续乘或加 \\
        \bottomrule
    \end{tabular}
\end{table}

构造 LL(1) 预测表时使用如下规则：

\begin{itemize}
\item 对产生式 $A \rightarrow \alpha$，若终结符 $a \in \mathrm{First}(\alpha)$，则把该产生式填入表项 $M[A][a]$；
\item 若 $\alpha$ 可以推导出空串，则对每个 $b \in \mathrm{Follow}(A)$，也把 $A \rightarrow \alpha$ 填入 $M[A][b]$。
\end{itemize}

如果某个表项被填入了两个不同的产生式，就说明文法不是 $LL(1)$，必须改写文法或使用 $LR$ 系解析器。

递归下降本身就是一种手工实现的 $LL(1)$ 解析：每个函数根据当前 Lookahead Token 选择进入哪个分支。正式的 $LL(1)$ 分析会用 $\mathrm{First}$ 集与 $\mathrm{Follow}$ 集计算预测表，把“该走哪个产生式”提前做成一张表。

% ---------------------------------------------------------------------
\subsection{LL(1) 解析器实现}

表驱动 $LL(1)$ 解析器需要两个栈：一个符号栈，一个输入 Token 栈（或指针）。伪代码如下：

\begin{lstlisting}
1. 开始时符号栈 = [expr, #]，输入指针指向第一个 Token；
2. 弹出符号栈顶 X：
   - 若 X 是终结符且等于当前 Token：消费输入；
   - 若 X 是非终结符：查预测表 M[X][当前Token]，
     得到产生式 X -> Y1 Y2 ... Yn，逆序压栈；
   - 若表项为空：语法错误。
3. 直到符号栈只剩 # 且输入结束。
\end{lstlisting}

例如 \texttt{factor} 面临当前 Token 是数字时，预测表会给出 \texttt{factor -> number}；当前 Token 是左括号时，给出 \texttt{factor -> '(' expr ')'}。这与递归下降中的 \texttt{switch(token)} 本质相同，只是把“判断逻辑”换成了查表。

本章完整示例采用递归下降实现，因为它更适合在单片机源码里逐步调试；如果脚本语言规模变大（增加 if、while 和函数调用），再迁移到由文法生成的 $LL(1)$ 或 $LR$ 解析器会更省事。

% ---------------------------------------------------------------------
\subsection{完整示例代码}

完整可运行的代码位于 \texttt{Chapter-script-engine/表达式脚本引擎/}：

\begin{itemize}
	\item \texttt{expr\_parser.h}：变量、错误码、公共 API；
	\item \texttt{expr\_parser.c}：词法扫描与递归下降解析；
	\item \texttt{main.c}：变量设置与多条表达式测试。
\end{itemize}

PC 上编译运行结果：

\begin{lstlisting}
	10 + 2 * (3 + x)             = 36
	(a + b) * 2                  = 24
	100 / 4 - 3 * 2              = 19
	1 + (2 + (3 + (4 + (5 + (6 + (7)))))) = 28
	10 / 0                       -> 除数为零
	x + y                        = 13
\end{lstlisting}

把读取表达式与输出结果的接口改成串口后，这套代码就能在单片机上提供“在线公式计算”能力；再结合第~\ref{ch:coroutine-scheduler}~章的协程模型，还可以让脚本计算在低优先级任务中执行，不阻塞其他控制逻辑。

% ---------------------------------------------------------------------
\section{解析器方案的选择}

在嵌入式环境中实现语法解析，通常有两种思路：手写递归下降解析器，或借助解析器生成工具（如 Yacc/Bison）生成表驱动解析器（典型为 $LALR(1)$）。表~\ref{tab:parser-choice} 从几个关键维度对这两种方案做了对比。

\begin{table}[htbp]
	\centering \small
	\caption{递归下降与 $LALR(1)$ 表驱动解析器对比}
	\label{tab:parser-choice}
	\begin{tabular}{lll}
		\toprule
		对比项 & 递归下降 & $LALR(1)$ 表驱动 \\
		\midrule
		解析方向 & 自顶向下展开 & 自底向上归约 \\
		推导方式 & 最左推导 & 最右推导的逆过程 \\
		解析模型 & 预测展开（可带回溯） & 移进-归约 \\
		左递归 & 必须消除 & 允许，无需改写文法 \\
		分析表 & 不需要，直接编码 & 需要 ACTION 与 GOTO 表 \\
		表体积 & 无 & 中等（可优化压缩） \\
		文法表达能力 & 较弱（需手工改写） & 强（涵盖大多数编程语言文法） \\
		资源占用（Flash/RAM） & 极小（仅代码） & 较大（表数据 + 运行时栈） \\
		错误信息 & 可定制，直观 & 通常生硬，难以定位 \\
		调试与维护 & 代码结构清晰，易于单步调试 & 表生成后黑盒，冲突排查困难 \\
		适用 MCU & 非常合适 & 需谨慎评估，大型表可能撑爆 Flash \\
		\bottomrule
	\end{tabular}
\end{table}

对于表达式这种小规模文法，递归下降足以胜任，且代码轻量、易读；而对于包含 if、while、函数调用等复杂结构的完整语言，$LALR(1)$ 生成器能减少大量手写工作量。但无论选哪一种，都需要清醒认识到：编译原理的“能力”不等于“可用性”，在实际 MCU 上，Flash 和 RAM 的限制往往比文法覆盖能力更关键。

以下是一些需要留意的权衡点：
\begin{enumerate}
	\item 递归下降要求文法无左递归，通常需改写文法并提取左公因子，这可能导致语法树结构与原始文法直觉不一致；
	\item $LALR(1)$ 虽然能力强大，但一旦出现归约-归约冲突，调试分析表十分困难，且生成的报错信息往往不友好，需要额外封装错误恢复；
	\item 无论哪种 $k=1$ 解析器，都只能向前看一个 Token，某些需要更长上下文的语言结构（如歧义型表达式）需结合词法层处理或提升 $k$ 值；
	\item 分析表只解决“语法对不对”，语义检查（类型、作用域、符号表）仍需单独编码实现，生成器不提供帮助。
\end{enumerate}

考虑到本章示例仅需解析简单表达式，且文法固定，手写递归下降比引入解析器生成器更节省 Flash，也能让读者直观地看到“每个函数对应一条文法规则”的结构，因此本章最终采用递归下降实现。关于 $\mathrm{First}$ 与 $\mathrm{Follow}$ 的构造算法以及更深入的编译器理论，可参考经典教材“龙书”：

\begin{itemize}
	\item 中文书名：《编译原理》（“龙书”中文翻译版，机械工业出版社）
	\item 英文原名：\textit{Compilers: Principles, Techniques, and Tools}， Second Edition
	\item 作者：Alfred V. Aho、Monica S. Lam、Ravi Sethi、Jeffrey D. Ullman
\end{itemize}

% ---------------------------------------------------------------------
\section{借助 Lua 框架实现}

当脚本功能从“计算表达式”升级为“完整控制逻辑”时，不建议继续手写解释器，而应使用成熟的 Lua：

\begin{itemize}
\item Lua 核心小：完整 Lua 5.x 可以裁剪到几十 KB Flash，适合资源较多的 32 位 MCU；
\item C API 清晰：固件通过 \texttt{lua\_register} 导出函数，脚本调用这些函数控制 LED、电机、传感器；
\item 安全性：Lua 自带语法错误、运行时错误与保护模式（\texttt{lua\_pcall}），不容易因为脚本异常导致裸机崩溃。
\item 内置完整前端：词法、文法、语法分析、字节码编译和虚拟机都是现成的，无需自己维护递归下降器；
\item 协程支持：Lua 原生协程适合“等待外设或定时器后继续执行”的脚本流程；
\item 生态成熟：错误信息、调试接口、内存统计等工具齐全，社区资料丰富。
\end{itemize}

与“自己维护一个表达式解析器”相比，Lua 的优势可以概括为表~\ref{tab:lua-advantage}~。

\begin{table}[htbp]
\centering \caption{自写表达式解释器与 Lua 的对比} \label{tab:lua-advantage}
    \begin{tabular}{lll}
        \toprule
        对比项 & 本章自写解释器 & Lua \\
        \midrule
        语法能力 & 四则运算表达式 & 完整脚本语言 \\
        解析器 & 自己维护递归下降 & 官方 Lua 词法、语法与编译器 \\
        运行方式 & 边解析边求值 & 编译为字节码后由 VM 执行 \\
        内存管理 & 手动管理 & 自动 GC（可配置或裁剪） \\
        错误处理 & 自己返回错误码 & pcall 与 error 保护模式 \\
        协程 & 无 & 原生支持 \\
        移植成本 & 很低 & 需要裁剪与移植内存接口 \\
        \bottomrule
    \end{tabular}
\end{table}

Lua 的代价是需要移植内存分配、提供系统时间，并严格限制脚本可访问的资源。其架构与本章示例一致：词法、语法和 Lua 虚拟机把源码编译成字节码，解释器执行字节码；导出给脚本的函数就是 C 侧的“系统调用”。

% ---------------------------------------------------------------------
\section{Python 解析生成器}

仓库里的 \texttt{Chapter-script-engine/parser.py} 是一个单文件工具：它负责加载 \texttt{.gram}、计算最终文法和预测分析表，并生成自带 Lexer 的非递归 C 解析器。生成出的单个 C 文件会把原 \texttt{variable\_runtime.c/h} 的完整实现内嵌，并附带一个演示 \texttt{main()}，因此不需要额外复制运行时源码即可直接编译运行。生成结果既可以继续接收调用者准备的 TOKEN 流，也可以通过 \texttt{Parser\_RunText} 直接解析脚本文本。

文法只使用 Yacc 风格但去掉 \texttt{|} 分隔符的 \texttt{.gram} 文本格式：

\begin{lstlisting}
%define PLUS "+"
%define MINUS "-"
%define STAR "*"
%define SLASH "/"
%define LPAREN "("
%define RPAREN ")"
%start expr

%var X int32 = 0
%var table int32[16] = { 1, 2, 3 }
%var pid struct {
    kp float
    ki float
    out int32
} = { kp: 1.0, out: 100 }
%var config dict[int32] = { mode: 1 }

expr := expr PLUS term      { $$ = $1 + $3; }
expr := expr MINUS term     { $$ = $1 - $3; }
expr := term                { $$ = $1; }

term := term STAR factor    { $$ = $1 * $3; }
term := factor              { $$ = $1; }

factor := NUMBER            { $$ = $1; }
factor := VARIABLE          { $$ = [$1]; }
factor := LPAREN expr RPAREN { $$ = $2; }
\end{lstlisting}

语法约定如下：

\begin{itemize}
\item \texttt{STRING}、\texttt{NUMBER} 与 \texttt{VARIABLE} 是默认终结符，不需要使用 \texttt{\%token} 重复声明；
\item \texttt{\%token}：声明只由外部 TOKEN 流提供的其他终结符；
\item \texttt{\%define}：声明固定字符串 Token，例如 \texttt{\%define EQ "==|="}；其中只支持竖线表示“或”，不支持正则表达式；
\item \texttt{\%start}：起始非终结符；
\item \texttt{\%tree}：让生成的 C 代码附带通用语法树 API（\texttt{Parser\_ParseTreeText} 等）；
\item \texttt{\%var}：声明脚本变量及其类型，可带初始值；
\item \texttt{\%var\_max\_count N}：设置生成代码中变量表的固定容量（默认 32）；
\item 每条产生式单独写成 \texttt{lhs := rhs \{ 动作 \}}，不使用 \texttt{|}，也不需要结尾 \texttt{;}；
\item \texttt{\%empty} 表示空产生式，例如 \texttt{else\_part := \%empty \{ \$\$ = int32(0); \}}。
\end{itemize}

\subsection{默认 Lexer 与动态数字}

\texttt{Parser\_RunText} 会自动识别双引号字符串、C 风格变量名和数字。整数默认创建为 \texttt{int32\_t}；整数后缀 \texttt{u} 创建为 \texttt{uint32\_t}；带小数点或科学计数法的数字默认创建为 \texttt{float}；后缀 \texttt{d} 则创建为 \texttt{double}。例如 \texttt{1}、\texttt{1u}、\texttt{1.2}、\texttt{1e3} 和 \texttt{1.2d} 会得到不同的 \texttt{VarType\_T}。正负号仍由文法作为运算符处理，以免把 \texttt{1-2} 的减号误收进数字。

多个固定字符串写在同一个 \texttt{\%define} 中时，Lexer 使用最长匹配。因此 \texttt{\%define EQ "==|="} 会先尝试 \texttt{==}，不会把它拆成两个 \texttt{=}。

\subsection{变量定义与固定类型赋值}

\begin{lstlisting}
def[$1] = $3;             /* 类型由 $3 推导 */
def[$1, int32] = $3;      /* 定义为 int32，并转换初值 */
def[$1, int32[8]] = $3;   /* 定义固定元素类型与长度的数组 */
def[$1, $2] = $3;         /* 复制 $2 的简单类型或数组类型 */
[$1] = $3;                /* 给已有变量赋值 */
$$ = [$1];                /* 读取变量作为产生式结果 */
delete[$1];               /* 删除普通变量 */
\end{lstlisting}

通过 \texttt{\%var} 预声明的变量属于系统默认变量：类型固定，不能被 \texttt{def} 重新定义，也不能用 \texttt{delete} 删除；脚本里尝试覆盖或删除系统变量会返回 \texttt{PARSER\_ERR\_SYSTEM}。运行时通过 \texttt{def} 创建的普通变量则可以再次 \texttt{def}，并且允许改变类型；\texttt{delete[\$1]} 会删除普通变量，之后可以重新定义。外部 C 代码也可以调用 \texttt{Parser\_DeleteVariable} 删除普通变量。\texttt{def[\$1, \$2]} 可以从已有的简单类型、数组、字典、结构体或嵌套复合值复制类型；显式字面类型写法仍用于简单类型和定长数组。

“一个非终结符的多条候选”不再用 \texttt{|} 连接，而是重复写 \texttt{lhs :=}：

\begin{lstlisting}
expr := expr PLUS term      { $$ = $1 + $3; }
expr := expr MINUS term     { $$ = $1 - $3; }
expr := term                { $$ = $1; }
\end{lstlisting}

编译器会直接认出“\texttt{lhs} 出现在自身 RHS 开头”的左递归，自动做以下展开：

\begin{enumerate}
\item 先解析一条“基产生式”（例如 \texttt{expr := term}），得到左值；
\item 循环检查当前 TOKEN 是否属于续产生式的开始符集合；
\item 每次命中续产生式时，把上一轮结果当作新的 \texttt{\$1}，解析后续操作数后重新执行动作；
\item 不再命中续产生式时结束折叠。
\end{enumerate}

Python 端会把直接左递归转换成带尾非终结符的最终文法。C 端不再生成 \texttt{Parse\_expr} 一类递归函数，而是把终结符、非终结符和语义动作标记压入显式分析栈。上一轮的结果作为尾规则的继承值继续折叠，因此 \texttt{a - b - c} 仍解释为 \texttt{(a-b)-c}。

\begin{lstlisting}[style=c]
/* 概念性伪代码：表驱动主循环 */
while (分析栈不为空) {
    item = 弹出分析栈;
    if (item 是终结符) 匹配当前 Token;
    else if (item 是非终结符) 按预测表展开产生式;
    else 执行产生式语义动作;
}
\end{lstlisting}

这种实现的 C 调用深度固定，解析深度只受 \texttt{PARSER\_MAX\_STACK\_DEPTH} 限制。生成器仍拒绝间接左递归和预测表冲突，因为本工具的目标是轻量 LL(1)，而不是制造 YACC。

\subsection{复杂变量与初始化列表}

\texttt{\%var} 支持标量、定长数组、结构体和字典。变量类型会直接生成对应的运行时构建函数，并使用 \texttt{VarValue\_T} 动态运行时保存：

\begin{lstlisting}
%var X int32 = 0                  // 标量
%var table int32[16] = {1,2,3}   // 数组，少则补默认值
%var grid int16[4][3]             // 二维数组
%var pid struct {
    kp float
    ki float
    out int32
} = { kp: 1.0, out: 100 }        // 结构体：按字段名初始化
%var config dict[int32] = { mode: 1 }  // 字典键值初始化
\end{lstlisting}

聚合初始化可以递归嵌套，例如“结构体数组”“数组里的结构体”“多维数组”都能直接写：

\begin{lstlisting}
%var team struct {
    id int32
    score float
}[2] = { {id:1, score:2.0}, {id:2, score:3.5} }

%var matrix int32[2][2] = { {1,2}, {3,4} }
\end{lstlisting}

\subsection{语义动作}

动作写在 \texttt{\{ ... \}} 里，支持 \texttt{\$\$}、\texttt{\$1、\$2...} 引用和完整的括号优先级：

\begin{lstlisting}
$$ = ($1 + $2) * $3 / ($1 - $2);   /* 自动展开为 Var_Add、Var_Mul 与 Var_Div */
$$ = $1.int32 + $2.int32;          /* 显式按 int32 做 C 标量运算 */
$1.int32 = 42;                     /* 类型化赋值 */
$$ = int32(0);                     /* 等价 Var_CreateInt(0) */
$$ = $1 ^ $2;                      /* 未显式类型：生成期告警 */
$$ = $1.int32 ^ $2.int32;          /* 显式 int32 位运算 */
\end{lstlisting}

各运算符到运行时的映射关系：

\begin{itemize}
\item \texttt{+ - * /} 依次对应 \texttt{Var\_Add}、\texttt{Var\_Sub}、\texttt{Var\_Mul} 与 \texttt{Var\_Div}；
\item \texttt{\%} 对应 \texttt{Var\_Mod}：整数执行取模，浮点执行 \texttt{fmod}；
\item \texttt{\&\& || !} 依次对应 \texttt{Var\_And}、\texttt{Var\_Or} 和 \texttt{Var\_Not}；
\item \texttt{[\$1]}：读取已有变量；放在赋值号左侧时执行固定类型转换后赋值；
\item \texttt{def[\$1, int32]}：定义具有固定类型的新变量，也可用 \texttt{\$n} 提供类型模板；
\item \texttt{\$1.int32}：按 int32 读取和写入运行时值的数值；
\item 未显式类型就对运行时值做位运算，生成时会告警；建议先写 \texttt{\$1.int32} 这样的类型字段，再参与异或、位与等运算。
\end{itemize}

生成结果包含 Token 枚举、预测分析表、显式分析栈、变量表等，以及 \texttt{Parser\_Run} 和 \texttt{Parser\_RunText} 两个入口。同一个 C 文件还会依次内嵌 \texttt{variable\_runtime.h} 与 \texttt{variable\_runtime.c} 的全部实现，并在末尾生成演示 \texttt{main()}：编译后直接运行，即可解析表达式并打印结果；若要把该 C 文件嵌入其他工程，可在 \texttt{\#include} 前定义 \texttt{PARSER\_SKIP\_DEMO\_MAIN} 来跳过演示入口。语义值使用 \texttt{VarValue\_T}（无 GC，由 \texttt{Var\_Free} 手动释放）。外部 C 程序可以用 \texttt{Var\_GetInt32}、\texttt{Var\_GetUint32}、\texttt{Var\_GetFloat}、\texttt{Var\_GetDouble} 等 Getter 做带范围检查的访问。

\subsection{生成命令}

生成默认四则运算文法：

\begin{lstlisting}[language=Python]
python Chapter-script-engine/parser_gen.py -o parser_gen.c
\end{lstlisting}

生成 \texttt{.gram} 文本文法：

\begin{lstlisting}[language=Python]
python Chapter-script-engine/parser_gen.py \
    --grammar-file Chapter-script-engine/grammar_calc.gram \
    --report parser_calc.report \
    -o parser_calc.c
\end{lstlisting}

生成出的 \texttt{parser\_calc.c} 已经自包含，可直接编译并运行演示程序：

\begin{lstlisting}[language=Python]
gcc parser_calc.c -lm -o parser_calc
./parser_calc "(1+2)*3"
\end{lstlisting}

生成器默认在控制台输出消除直接左递归后的最终产生式、FIRST 集合、FOLLOW 集合等信息；\texttt{--report} 会把同样内容保存为 UTF-8 文本，便于读者逐项核对。

完整示例见 \texttt{Chapter-script-engine/grammar\_calc.gram}；包含 IF/FOR、布尔与比较运算的高级示例见 \texttt{Chapter-script-engine/grammar\_advanced.gram}。两类文法都通过 \texttt{--grammar-file} 交给生成器，统一使用 \texttt{.gram} 文本格式。

\begin{anybox}{为什么不把运行时与生成器展开讲解}
	细心的读者可能已经注意到，本章的 \texttt{parser.py} 已经不再只是“文法转 C”的翻译器：它把 \texttt{variable\_runtime.c/h} 的实现收进自身，生成的 C 文件又同时携带动态类型运行时和演示 \texttt{main()}。正文中仍然只浅尝辄止地介绍它们的存在，这并不是遗漏，而是刻意的边界划定。
	
	\texttt{VarValue\_T} 运行时实现了一套完整的动态类型系统——支持多种类型、深拷贝、切片拼接、类型自动提升等。这套代码本身已经构成了一个“微型语言运行时”，其设计涉及内存所有权、容器扩容、类型转换等话题，如果要讲透，至少需要再写半本书。本章的核心目标是让读者理解“如何用递归下降解析表达式”，而不是“如何设计一个动态语言虚拟机”，因此运行时的细节被封装进生成器和生成产物，作为“愿意深究的读者可以翻阅”的材料。
	
	\texttt{parser.py} 则是“工具链的工具链”——它把人类可读的文法描述翻译成带 Lexer、显式分析栈、内嵌运行时和演示 \texttt{main()} 的 C 代码。这套生成器涉及 FIRST/FOLLOW 集合、预测表、左递归转换、动作表达式翻译等编译原理话题，每一部分都足以单独成章。如果把这些内容全部铺进正文，这章就会从“单片机脚本引擎”变成“编译原理实践教程”，彻底脱离本书“嵌入式 C 生存实战”的定位。
	
	因此，作者的选择是：在正文中保留生成器的“使用方式”和“设计思路”，让读者知道存在这样一套工具，但把实现细节留在源码里。工具的价值在于被使用，而不是被背诵。对于有兴趣深挖的读者，仓库里的源码和注释是比任何书本都更直接的老师。
	
	这种“正文讲思想，仓库供细节”的做法，也是全书一以贯之的风格：只把必须由作者推导给你看的核心逻辑（如状态机、内存池、协程切换）写进纸面，把“查手册就能得到”的实现细节留给代码本身。
\end{anybox}

% ---------------------------------------------------------------------
\section{解析器的选择与工程边界}

如果你读到这里觉得“这章怎么突然就停了”，没错，是我故意停的。

因为当我把递归下降、$LL(1)$ 以及 First 与 Follow 集合全部铺完之后，我必须面对一个诚实的问题：在单片机上，你到底该不该手写解析器？答案分三种情况，如表~\ref{tab:parser-decision}~所示。

\begin{table}[htbp]
	\centering
	\small
	\caption{单片机解析方案选择}
	\label{tab:parser-decision}
	\begin{tabularx}{\textwidth}{l X l}   % X列自动换行
		\toprule
		你的需求 & 推荐方案 & 理由 \\
		\midrule
		只需要算四则运算 + 简单变量 & 用本章所讲述的Python的文法生成器生成 & 代码量小、无依赖、实时性可控 \\
		需要完整脚本 & 直接使用 Lua 或 MicroPython & 自己写的虚拟机很难超过成熟 Lua \\
		解析行业标准 DSL，硬实时要求高 & Python生成器 + 文法描述生成 C 解析器 & 编译期展开、零动态内存 \\
		\bottomrule
	\end{tabularx}
\end{table}

\begin{anybox}{那龙书是干什么用的？}
龙书不是用来“照着实现”的，它是用来查文法怎么写的。当你写的文法出现左递归消除不掉、二义性搞不定时，翻龙书的第 4 章（LL）。它不是教材，它是字典。
\end{anybox}

\begin{anybox}{那 Python 生成器是干什么用的？}
它是拿来把你写的文法“翻译”成 C 代码的。你写好生成式，它帮你生成\texttt{ParseExpr}、\texttt{ParseTerm}、\texttt{ParseFactor}，附带行号追踪和错误恢复骨架。你只需要维护文法文件，不需要手改 C 代码。
\end{anybox}

\begin{anybox}{为什么硬实时 DSL 不适合直接套用 Lua？}
因为 G 代码的宏展开（例如 \texttt{\#1 = 10; G01 X[\#1 + 5]}）必须在解析期完成展开，然后把展开后的结果直接变成插补器的输入。Lua 的运行时变量查表会在每次插补周期引入不可预测的延迟，而生成器在编译期就把 \texttt{\#1} 变成 \texttt{axis\_pos[1]} 的直接数组访问——少了整整一个“解释执行”的开销层级。
\end{anybox}

仓库中 \texttt{Chapter-script-engine/parser.py} 的更多实现细节见附录~\ref{app:parser-generator-runtime}。

% ---------------------------------------------------------------------
\begin{problemset}
\item 解释终结符、非终结符、产生式和开始符号四个概念，并用一个简单加法表达式写出对应的文法。
\item 递归下降解析器为什么要把乘法放在比加法更深的“函数层级”中？如果两个运算符优先级相同，如何保证结合性？
\item 直接左递归文法（如 \texttt{expr := expr PLUS term}）在递归下降中为什么不能直接实现？本章的生成器如何把它变成“基产生式 + 循环”？
\item First 与 Follow 集合分别用于解决什么问题？LL(1) 冲突出现时，通常有哪几种处理手段？
\item .gram 文法中，左递归、\texttt{\%var} 复杂变量和 \texttt{int32(...)} 动作各让哪些原本需要手工维护的中间表示或动作代码得以省略？
\end{problemset}
